% ============================================================
%  Divisibilidade — Layout Hero | PratiqueJa
%  Engine: XeLaTeX  |  Font: Inter via fontspec
% ============================================================
\documentclass[12pt]{article}

\usepackage[brazil]{babel}

\usepackage{fontspec}
\setmainfont[
  BoldFont       = Inter-Bold,
  ItalicFont     = Inter-Italic,
  BoldItalicFont = Inter-BoldItalic,
]{Inter-Regular}
\usepackage{microtype}

% ── Geometria: margens visuais equilibradas ──────────────────
\usepackage[
  a4paper,
  top=1.5cm,
  bottom=2.8cm,
  left=1.5cm,
  right=1.5cm,
]{geometry}
\setlength{\footskip}{1.3cm}   % rodapé a ~1.5 cm da borda inferior

% ── Cores ────────────────────────────────────────────────────
\usepackage[dvipsnames,svgnames,table]{xcolor}

\definecolor{primary}{HTML}{2563EB}
\definecolor{primarydark}{HTML}{1E40AF}
\definecolor{heroback}{HTML}{1D4ED8}
\definecolor{lightblue}{HTML}{60A5FA}
\definecolor{navybrand}{HTML}{051C4E}
\definecolor{softbg}{HTML}{F0F5FF}
\definecolor{chipbg}{HTML}{EFF6FF}
\definecolor{chipborder}{HTML}{BFDBFE}
\definecolor{bodytext}{HTML}{374151}
\definecolor{titletext}{HTML}{07142F}
\definecolor{mutedtext}{HTML}{9CA3AF}
\definecolor{exborder}{HTML}{E5E7EB}
\definecolor{exbg}{HTML}{F8FAFC}
\definecolor{badgepurple}{HTML}{7C3AED}
\definecolor{badgegreen}{HTML}{059669}

\colorlet{iris}{primary}
\colorlet{bluebell}{lightblue}

% ── Matemática ───────────────────────────────────────────────
\usepackage{amsmath,amssymb}

% ── Layout ───────────────────────────────────────────────────
\usepackage{multicol}
\usepackage{setspace}
\usepackage{parskip}
\usepackage{graphicx}
\usepackage{array}

% ── Caixas estilizadas ───────────────────────────────────────
\usepackage[most]{tcolorbox}
\tcbuselibrary{skins, breakable}

% ── Rodapé ───────────────────────────────────────────────────
\usepackage{fancyhdr}
\pagestyle{fancy}
\fancyhf{}
\renewcommand{\headrulewidth}{0pt}
\renewcommand{\footrulewidth}{0.6pt}
\fancyfoot[L]{\small\color{mutedtext}© Copyright, Pratique Já}
\fancyfoot[C]{\color{primary}\textbf{pratiqueja.com}}

% ============================================================
%  COMANDOS PERSONALIZADOS
% ============================================================

% ── Logo ─────────────────────────────────────────────────────
\newcommand{\brandlogo}{%
  {\fontsize{22}{26}\selectfont\bfseries
    \textcolor{navybrand}{Pratique}\textcolor{primary}{Já}}%
}

% ── Hero header — fundo branco com bordas azuis ──────────────
\newcommand{\heroheader}[2]{%
  \begin{tcolorbox}[
    enhanced, arc=0pt, boxrule=0pt,
    colback=white, colframe=white,
    left=18pt, right=18pt, top=12pt, bottom=12pt,
    borderline north={2pt}{0pt}{primary},
    borderline south={3pt}{0pt}{primary},
  ]
    \begin{minipage}[c]{0.26\linewidth}
      \centering\brandlogo
    \end{minipage}%
    \hfill
    \begin{minipage}[c]{0.70\linewidth}
      \centering
      {\fontsize{30}{36}\selectfont\bfseries\color{heroback}#1}\par
      \vspace{6pt}
      {\color{primary}\rule{0.5\linewidth}{1.5pt}}\par
      \vspace{6pt}
      {\color{bodytext}\large\itshape #2}
    \end{minipage}
  \end{tcolorbox}%
  \vspace{8pt}%
}

% ── Badge inline — equivalente ao .tr-badge ──────────────────
\newcommand{\badge}[2]{%  {texto}{cor}
  \tcbox[
    on line, enhanced, arc=6pt, boxrule=0pt,
    colback=#2, colframe=#2,
    left=4pt, right=4pt, top=3pt, bottom=3pt,
    nobeforeafter,
  ]{\footnotesize\bfseries\color{white}#1}%
}

% ── Cabeçalho de seção — equivalente ao .tr-section-header ───
% #1=badge  #2=badge-color  #3=título  #4=subtítulo
\newcommand{\TW}[4]{%
  \vspace{10pt}%
  \noindent\badge{#1}{#2}\hspace{6pt}%
  {\color{titletext}\bfseries\large #3}\par
  \noindent\hspace{2pt}{\color{mutedtext}\small #4}\par
  \vspace{4pt}%
  {\color{exborder}\hrule height 1.5pt}%
  \vspace{6pt}%
}

% ── Rótulo interno h4 ────────────────────────────────────────
\newcommand{\SH}[1]{%
  \vspace{7pt}%
  \noindent{\color{mutedtext}\fontsize{8}{10}\selectfont\bfseries #1}\par
  \vspace{3pt}%
}

% ── Caixa de regra — equivalente ao .tr-rule ─────────────────
\newcommand{\TM}[1]{%
  \begin{tcolorbox}[
    enhanced, arc=0pt, boxrule=0pt,
    colback=softbg, colframe=softbg,
    left=10pt, right=8pt, top=8pt, bottom=8pt,
    borderline west={4pt}{0pt}{primary},
  ]
    \color{primarydark}\small #1
  \end{tcolorbox}%
  \vspace{2pt}%
}

% ── Caixa de exemplo — equivalente ao .tr-ex ─────────────────
\newcommand{\trex}[1]{%
  \vspace{2pt}%
  \begin{tcolorbox}[
    enhanced, arc=6pt, boxrule=0.5pt,
    colback=exbg, colframe=exborder,
    left=8pt, right=8pt, top=5pt, bottom=5pt,
  ]
    \small\color{bodytext}#1
  \end{tcolorbox}%
  \vspace{1pt}%
}

% ── Chip de divisor — equivalente ao .tr-divisor ─────────────
\newcommand{\divisorchip}[1]{%
  \tcbox[
    on line, enhanced, arc=6pt, boxrule=0.5pt,
    colback=chipbg, colframe=chipborder,
    left=5pt, right=5pt, top=2pt, bottom=2pt,
    nobeforeafter,
  ]{\small\bfseries\color{primarydark}#1}\hspace{3pt}%
}

% ============================================================
\begin{document}
\pagenumbering{gobble}
\setstretch{1.3}
\color{bodytext}

\heroheader{Divisibilidade}{Regras, Divisores e Aplicações}

\setlength{\columnsep}{22pt}
\setlength{\columnseprule}{0.4pt}
\def\columnseprulecolor{\color{chipborder}}

\begin{multicols}{2}

%─────────────────────────────────────────────────────────────
\TW{Def}{badgepurple}{Definição}{O que é divisibilidade?}

\TM{Um número é divisível por um determinado divisor se e somente se o
resto da divisão for igual a zero.}

Dizemos que $a$ é divisível por $b$ quando existe um inteiro $q$ tal
que $a = b \cdot q$. Nesse caso, $b$ é chamado de \textbf{divisor} de
$a$ e $a$ é chamado de \textbf{múltiplo} de $b$.

%─────────────────────────────────────────────────────────────
\TW{2}{primary}{Divisibilidade por 2}{Números pares}

\TM{Todo número \textbf{par} é divisível por 2 — quando o último
dígito é 0, 2, 4, 6 ou 8.}

\SH{Exemplos}

\trex{$2,\ 4,\ 6,\ 8,\ 10,\ \ldots$}
\trex{$13\textcolor{primary}{8}$ → último dígito 8 → divisível ✓}
\trex{$47\textcolor{primary}{5}$ → último dígito 5 → não divisível ✗}

%─────────────────────────────────────────────────────────────
\TW{3}{primary}{Divisibilidade por 3}{Soma dos algarismos}

\TM{Um número é divisível por 3 quando a \textbf{soma dos seus
algarismos} for múltipla de 3.}

\SH{Exemplos}

\trex{$132 \Rightarrow 1+3+2 = \textcolor{primary}{6}$ → múltiplo de 3 ✓}
\trex{$531 \Rightarrow 5+3+1 = \textcolor{primary}{9}$ → múltiplo de 3 ✓}
\trex{$124 \Rightarrow 1+2+4 = \textcolor{primary}{7}$ → não múltiplo de 3 ✗}

%─────────────────────────────────────────────────────────────
\TW{4}{primary}{Divisibilidade por 4}{Dois últimos dígitos}

\TM{Um número é divisível por 4 quando o número formado pelos
\textbf{2 últimos dígitos} é divisível por 4.}

\SH{Exemplos}

\trex{$5\textcolor{primary}{24}$ → 24 é múltiplo de 4 ✓}
\trex{$13\textcolor{primary}{36}$ → 36 é múltiplo de 4 ✓}
\trex{$5\textcolor{primary}{14}$ → 14 não é múltiplo de 4 ✗}

%─────────────────────────────────────────────────────────────
\TW{5}{primary}{Divisibilidade por 5}{Último dígito 0 ou 5}

\TM{Um número é divisível por 5 quando o \textbf{último dígito} é
igual a \textbf{0} ou \textbf{5}.}

\SH{Exemplos}

\trex{$32\textcolor{primary}{5}$ → termina em 5 ✓}
\trex{$14\textcolor{primary}{0}$ → termina em 0 ✓}
\trex{$78\textcolor{primary}{3}$ → termina em 3 ✗}

%─────────────────────────────────────────────────────────────
\TW{6}{primary}{Divisibilidade por 6}{Divisível por 2 e por 3}

\TM{Um número é divisível por 6 quando é \textbf{divisível por 2}
(número par) \textbf{e por 3} (soma dos algarismos múltipla de 3)
simultaneamente.}

\SH{Exemplos}

\trex{$654 \Rightarrow 6+5+4=\textcolor{primary}{15}$ — mult.\ de 3 e par ✓}
\trex{$312 \Rightarrow 3+1+2=\textcolor{primary}{6}$ — mult.\ de 3 e par ✓}
\trex{125 → ímpar → não divisível por 6 ✗}

%─────────────────────────────────────────────────────────────
\TW{7}{primary}{Divisibilidade por 7}{Regra do duplo}

\TM{Multiplique o \textbf{último dígito por 2} e subtraia do número
restante (sem o último dígito). Repita até que o resultado seja
visivelmente divisível ou não por 7.}

\SH{Exemplos}

\trex{$16\textcolor{primary}{1} \Rightarrow 1{\cdot}2=2
\Rightarrow 16-2=\textcolor{primary}{14}$ → múltiplo de 7 ✓}
\trex{$109\textcolor{primary}{9} \Rightarrow 9{\cdot}2=18
\Rightarrow 109{-}18=91$\\
→ repete: $9\textcolor{primary}{1}
\Rightarrow 9{-}2=\textcolor{primary}{7}$ → múltiplo de 7 ✓}

%─────────────────────────────────────────────────────────────
\TW{8}{primary}{Divisibilidade por 8}{Três últimos dígitos}

\TM{Um número é divisível por 8 quando o número formado pelos
\textbf{3 últimos dígitos} é divisível por 8.}

\SH{Exemplos}

\trex{$7\textcolor{primary}{128}$ → 128 é múltiplo de 8 ✓}
\trex{$5\textcolor{primary}{256}$ → 256 é múltiplo de 8 ✓}
\trex{$3\textcolor{primary}{100}$ → 100 não é múltiplo de 8 ✗}

%─────────────────────────────────────────────────────────────
\TW{9}{primary}{Divisibilidade por 9}{Soma dos algarismos}

\TM{Um número é divisível por 9 quando a \textbf{soma dos seus
algarismos} for múltipla de 9.}

\SH{Exemplos}

\trex{$432 \Rightarrow 4+3+2=\textcolor{primary}{9}$ → múltiplo de 9 ✓}
\trex{$6561 \Rightarrow 6+5+6+1=\textcolor{primary}{18}$ → múltiplo de 9 ✓}
\trex{$124 \Rightarrow 1+2+4=\textcolor{primary}{7}$ → não múltiplo de 9 ✗}

%─────────────────────────────────────────────────────────────
\TW{10}{primary}{Divisibilidade por 10}{Último dígito 0}

\TM{Um número é divisível por 10 quando o \textbf{último dígito é 0}.}

Como $10 = 2 \times 5$, ser divisível por 10 equivale a ser divisível
por 2 \emph{e} por 5 simultaneamente.

\SH{Exemplos}

\trex{$15\textcolor{primary}{0}$ → termina em 0 ✓}
\trex{$487\textcolor{primary}{0}$ → termina em 0 ✓}
\trex{$32\textcolor{primary}{5}$ → termina em 5, divisível por 5 mas não por 10 ✗}

%─────────────────────────────────────────────────────────────
\TW{11}{badgegreen}{Divisibilidade por 11}{Soma alternada dos algarismos}

\TM{Um número é divisível por 11 quando a \textbf{soma alternada dos
seus algarismos} (alternando $+$ e $-$ da esquerda para a direita)
for múltipla de 11 (incluindo zero).}

\SH{Exemplos}

\trex{$253 \Rightarrow +2-5+3 = \textcolor{primary}{0}$ → múltiplo de 11 ✓}
\trex{$1452 \Rightarrow +1{-}4+5{-}2 = \textcolor{primary}{0}$ → múltiplo de 11 ✓}
\trex{$123 \Rightarrow +1{-}2+3 = \textcolor{primary}{2}$ → não múltiplo de 11 ✗}

%─────────────────────────────────────────────────────────────
\TW{Div}{badgegreen}{Divisores de um Número}{Como encontrar todos os divisores naturais}

Para encontrar todos os divisores, decomponha o número em fatores
primos. Para cada fator, multiplique \textbf{todos os divisores já
encontrados} por esse fator e acrescente os resultados novos.

\SH{Exemplo — divisores de 60}

\textbf{Passo 1} — Decomposição em fatores primos:

\vspace{4pt}
\noindent
$\begin{array}[t]{r|l}
60 & 2 \\ 30 & 2 \\ 15 & 3 \\ 5 & 5 \\
\hline 1 & \\
\end{array}$
\quad $60 = 2^2 \cdot 3 \cdot 5$

\vspace{6pt}
\textbf{Passo 2} — Coluna de divisores (inicia com 1, aplica cada fator):

\vspace{4pt}
\noindent
$\begin{array}[t]{r|l|l}
  & & \textcolor{primary}{1} \\
60 & \textcolor{primary}{2} & \textcolor{primary}{2} \\
30 & \textcolor{primary}{2} & \textcolor{primary}{4} \\
15 & \textcolor{primary}{3} & \textcolor{primary}{3,\ 6,\ 12} \\
 5 & \textcolor{primary}{5} & \textcolor{primary}{5,\ 10,\ 20,\ 15,\ 30,\ 60} \\
\hline
 1 &   & \\
\end{array}$

\SH{Todos os divisores de 60}

\noindent\divisorchip{1}\divisorchip{2}\divisorchip{3}\divisorchip{4}%
\divisorchip{5}\divisorchip{6}\divisorchip{10}\divisorchip{12}%
\divisorchip{15}\divisorchip{20}\divisorchip{30}\divisorchip{60}

\SH{Quantidade de divisores}

Se $n = p_1^{a_1} \cdot p_2^{a_2} \cdots$, o total de divisores é:

\TM{\centering$d(n) = (a_1+1)(a_2+1)\cdots$}

\trex{$60 = 2^2{\cdot}3^1{\cdot}5^1 \Rightarrow d(60)=(3)(2)(2)=\textcolor{primary}{12}$ divisores}
\trex{$36 = 2^2{\cdot}3^2 \Rightarrow d(36)=(3)(3)=\textcolor{primary}{9}$ divisores}

\end{multicols}

\end{document}
